\section{Modell und Datengrundlage}

\subsection{Architektur von Pythia-410M}

Untersucht wird das Modell \texttt{EleutherAI/pythia-410m} \cite{biderman2023pythia}. Es beruht auf der GPT-NeoX-Architektur und gehört zur Klasse der Decoder-Only-Transformer. Anders als die ursprüngliche \gls{transformer}-Architektur mit Encoder und Decoder verwendet GPT-NeoX ausschließlich Decoder-Blöcke mit maskierter Self-\gls{attention} \cite{vaswani2017attention}. Dadurch arbeitet das Modell autoregressiv und sagt auf Basis aller vorherigen \glspl{token} das jeweils nächste Token vorher. Tabelle~\ref{tab:arch} fasst die zentralen Architekturwerte zusammen.

\begin{table}[htbp]
  \centering
  \caption{Zentrale Architekturwerte von Pythia-410M.}
  \label{tab:arch}
  \begin{tabular}{ll}
    \toprule
    Eigenschaft & Wert \\
    \midrule
    Architekturtyp & Decoder-Only (GPT-NeoX) \\
    Vokabulargröße & 50.304 Token \\
    Hidden Size & 1.024 \\
    Transformer-Schichten & 24 \\
    \glspl{attentionhead} je Schicht & 16 \\
    Dimension je Kopf & 64 \\
    MLP-Zwischengröße & 4.096 \\
    Aktivierungsfunktion & GELU \\
    Positionskodierung & Rotary Positional Embeddings (RoPE) \\
    Maximale Kontextlänge & 2.048 Token \\
    \bottomrule
  \end{tabular}
\end{table}

Jedes Eingabetoken wird zunächst in einen 1.024-dimensionalen Embedding-Vektor projiziert. Jede der 24 Schichten enthält einen Multi-Head-Self-\gls{attention}-Block mit 16 Köpfen sowie ein nachgelagertes Feed-Forward-Netzwerk. Im Feed-Forward-Netzwerk wird die Repräsentation zunächst von 1.024 auf 4.096 Dimensionen erweitert und anschließend wieder auf 1.024 reduziert, wobei zwischen den beiden Projektionen die nichtlineare GELU-Funktion wirkt. Positionsinformationen werden über Rotary Positional Embeddings direkt in die Query- und Key-Repräsentationen eingebracht. Nach der letzten Schicht werden die Hidden States auf 50.304 Logits projiziert und über die Softmax-Funktion in Wahrscheinlichkeiten überführt.

\subsection{Parameterverteilung}

Eine Zählung der Parameter zeigt, wie sich die Kapazität des Modells auf die Komponenten verteilt. Insgesamt besitzt das Modell rund 405 Millionen Parameter. Tabelle~\ref{tab:params} ordnet diese den Hauptbestandteilen zu.

\begin{table}[htbp]
  \centering
  \caption{Verteilung der Parameter auf die Hauptkomponenten von Pythia-410M.}
  \label{tab:params}
  \begin{tabular}{lr}
    \toprule
    Komponente & Anzahl Parameter \\
    \midrule
    Embedding-Matrix (50.304 $\times$ 1.024) & 51.511.296 \\
    Transformer-Schichten (24 Schichten) & 302.309.376 \\
    Finale Layer-Normalisierung & 2.048 \\
    Unembedding-Matrix (1.024 $\times$ 50.304) & 51.511.296 \\
    \midrule
    Summe & 405.334.016 \\
    \bottomrule
  \end{tabular}
\end{table}

Innerhalb einer einzelnen \gls{transformer}-Schicht entfällt nur etwa ein Drittel der Parameter auf den \gls{attention}-Mechanismus, während rund zwei Drittel im Feed-Forward-Netzwerk liegen. Der größte Teil der Modellkapazität befindet sich somit nicht in der \gls{attention} selbst, sondern in den Feed-Forward-Blöcken, die wesentlich zur Speicherung und Verarbeitung von Wissen beitragen. Abbildung~\ref{fig:paramverteilung} stellt diese Verteilung grafisch dar.

\begin{figure}[htbp]
  \centering
  \includegraphics[width=0.75\textwidth]{parameterverteilung.png}
  \caption{Verteilung der Parameter von Pythia-410M auf Embedding, Unembedding, Attention und MLP.}
  \label{fig:paramverteilung}
\end{figure}

\subsection{Datensätze}

Für die Analysen werden drei Datensätze verwendet, die sich gegenseitig ergänzen.

\paragraph{Hu-\&-Liu-Opinion-Lexikon.}
Das Opinion-Lexikon von Hu und Liu liefert die positiven und negativen Wortlisten \cite{hu2004mining}. Da die Embedding-Analysen eine eindeutige Zuordnung von Wort zu Vektor erfordern, werden nur Wörter behalten, die der Tokenizer von Pythia-410M als genau ein Token darstellt. Tabelle~\ref{tab:hulius} zeigt die nach diesem Ein-Token-Filter verbleibenden Wörter. Diese Wörter bilden die Grundlage für die Embedding-Analysen, die lineare Probe und die logitbasierten Sentiment-Scores.

\begin{table}[htbp]
  \centering
  \caption{Anzahl der Hu-\&-Liu-Wörter nach dem Ein-Token-Filter für Pythia-410M.}
  \label{tab:hulius}
  \begin{tabular}{lr}
    \toprule
    Sentiment-Klasse & Anzahl Ein-Token-Wörter \\
    \midrule
    positiv & 865 \\
    negativ & 1.389 \\
    \midrule
    gesamt & 2.254 \\
    \bottomrule
  \end{tabular}
\end{table}

\paragraph{Sentiment Challenge Dataset.}
Der Challenge-Datensatz „Sentiment Analysis Is Not Solved“ enthält gezielt schwierige Sätze und annotiert jeden Satz mit sprachlichen Kategorien wie Negation, Ironie, Idiomen, Emoji oder gemischtem Sentiment \cite{barnes2019sentiment}. Die ursprünglichen Labels reichen von 0 für stark negativ bis 5 für stark positiv. Für die binäre Sentiment-Analyse werden neutrale Einträge entfernt und die verbleibenden Sätze nur noch als positiv oder negativ geführt. Dieser Datensatz dient der verhaltensbasierten Bewertung in Kapitel 4.

\paragraph{Counterfactually Augmented Data.}
Der CAD-Datensatz enthält Paare aus einem positiven und einem negativen Text, die sich nur durch wenige sentimenttragende Wörter unterscheiden \cite{kaushik2020learning}. Durch den nahezu identischen Kontext eignet er sich besonders gut, um Sentiment isoliert zu untersuchen. Der verwendete Teil umfasst 1.707 solcher gegenfaktischen Paare. Dieser Datensatz wird für die Aggregation der Logit-Lens-Scores in Kapitel 6 sowie für die Attention-Analyse in Kapitel 7 genutzt.

\subsection{Eingesetzte Werkzeuge}

Die Implementierung beruht auf PyTorch für Modellinferenz, Forward Hooks und Tensoroperationen sowie auf der Transformers-Bibliothek von HuggingFace für das Laden von Modell und Tokenizer. Das Hu-\&-Liu-Lexikon wird über NLTK bezogen. Die lineare Probe nutzt die logistische Regression aus scikit-learn. Die Abbildungen werden mit Matplotlib erstellt, und das Experiment-Tracking erfolgt über MLflow. Eine begleitende Dashboard-Anwendung macht ausgewählte Ergebnisse interaktiv zugänglich. Spezialbibliotheken für mechanistische Interpretierbarkeit werden bewusst nicht verwendet.
